\section{ĐẶT VẤN ĐỀ}

Tổng hợp tiếng nói từ văn bản (Text-to-Speech, viết tắt là TTS) là một trong những bài toán nền tảng thuộc lĩnh vực xử lý ngôn ngữ tự nhiên (NLP) và xử lý tín hiệu đa phương tiện, nhằm chuyển đổi tự động một thông điệp văn bản dưới dạng ký tự số sang tín hiệu sóng âm thanh tương ứng có thể nghe và hiểu được (Dutoit, 1997). Trong kỷ nguyên số hóa và trí tuệ nhân tạo hiện nay, công nghệ TTS giữ vai trò cốt lõi trong việc xây dựng giao diện tương tác người -- máy bằng giọng nói (Voice User Interface - VUI), hỗ trợ đắc lực cho người khiếm thị tiếp cận tri thức số, tích hợp vào các hệ thống tổng đài phản hồi thông minh (IVR), đọc sách điện tử và tin tức báo chí tự động, cũng như tạo giọng nói tương tác cho các trợ lý ảo thông minh (Taylor, 2009).

Trong khuôn khổ học phần Cơ sở dữ liệu đa phương tiện, dữ liệu âm thanh (Audio) và văn bản (Text) đại diện cho hai đối tượng dữ liệu phi cấu trúc và bán cấu trúc tiêu biểu (Subrahmanian, 1998). Một hệ thống TTS toàn diện không chỉ đơn thuần giải quyết bài toán phát âm mà còn là bài toán thiết kế hệ thống quản trị cơ sở dữ liệu đa phương tiện: từ khâu tổ chức lưu trữ ngữ liệu song song (Parallel Audio-Text Corpus), trích xuất và lập chỉ mục các đặc trưng âm học số (Acoustic Feature Extraction and Indexing), cho đến việc tối ưu hóa truy vấn thời gian thực nhằm phục vụ tổng hợp tín hiệu mượt mà.

Đối với ngôn ngữ tiếng Việt, việc xây dựng hệ thống TTS phải đối mặt với nhiều thách thức mang tính chất âm học và ngôn ngữ học đặc thù:
\begin{enumerate}
    \item \textbf{Hệ thống thanh điệu phức tạp}: Tiếng Việt là một ngôn ngữ đơn lập có thanh điệu (Tonal language), bao gồm 6 thanh điệu cơ bản: Ngang (không dấu), Sắc, Huyền, Hỏi, Ngã và Nặng (Nguyễn Hữu Quỳnh, 2001). Thanh điệu trong tiếng Việt đóng vai trò khu biệt ngữ nghĩa sâu sắc: cùng một tổ hợp âm vị (ví dụ: \textit{ma, má, mà, mả, mã, mạ}) nhưng khi gắn với các thanh điệu khác nhau sẽ tạo thành các từ hoàn toàn khác biệt. Do đó, việc mã hóa chính xác thanh điệu là điều kiện tiên quyết để hệ thống tra cứu đúng bản ghi âm thanh tương ứng trong cơ sở dữ liệu.
    \item \textbf{Hiện tượng gián đoạn âm học tại ranh giới ghép nối}: Trong phương pháp tổng hợp ghép nối (Concatenative Speech Synthesis), các mẩu âm thanh độc lập được thu âm riêng rẽ khi ghép liên tiếp vào nhau sẽ xuất hiện sự gián đoạn đột ngột về pha tín hiệu (Phase discontinuity) và độ chênh lệch mức năng lượng (RMS energy mismatch) (Hunt và Black, 1996). Sự nhảy vọt biên độ tức thời này dẫn đến hiện tượng phát sinh tiếng nổ lách tách (click) hoặc tiếng bụp (pop), làm giảm sút nghiêm trọng chất lượng trải nghiệm của người nghe.
    \item \textbf{Ngữ điệu câu và hiện tượng đồng cấu âm (Coarticulation)}: Trong lời nói tự nhiên liên tục, các âm vị kế cận có xu hướng ảnh hưởng lẫn nhau do quán tính chuyển động liên tục của các bộ phận trong khoang miệng (môi, lưỡi, thanh môn) (Rabiner và Schafer, 2010). Việc ghép nối các đơn vị từ độc lập thường khiến ngữ điệu câu có cảm giác rời rạc, thiếu sự mượt mà và biến điệu tự nhiên của ngữ điệu tiếng Việt.
\end{enumerate}

Về phương pháp tổng hợp, lịch sử phát triển của công nghệ TTS đã ghi nhận ba hướng tiếp cận chủ đạo:
\begin{itemize}
    \item \textit{Tổng hợp tham số hoặc mô phỏng bộ máy phát âm (Formant/Articulatory Synthesis)}: Điển hình bởi hệ thống Klatt synthesizer hoặc phương pháp LPC (Linear Predictive Coding). Phương pháp này mô hình hóa đường truyền âm thanh của con người bằng các tham số toán học (Dutoit, 1997). Mặc dù tốc độ sinh nhanh và dung lượng lưu trữ cực nhỏ, giọng nói tổng hợp thường mang tính chất máy móc (robotic), khô cứng và thiếu đi tính tự nhiên.
    \item \textit{Tổng hợp ghép nối (Concatenative Speech Synthesis)}: Tiếp cận bằng cách lưu trữ một kho ngữ liệu các đơn vị âm thanh được thu âm thực tế (như từ, âm tiết, diphone hoặc triphone), sau đó lựa chọn và ghép nối tuần tự các đoạn âm thanh này để tạo thành câu hoàn chỉnh (Taylor, 2009). Điểm mạnh vượt trội của phương pháp này là bảo toàn 100\% chất giọng thực tế của người nói, không làm suy giảm chất lượng qua vocoder nhân tạo, đồng thời không đòi hỏi hạ tầng tính toán GPU phức tạp.
    \item \textit{Tổng hợp dựa trên mạng nơ-ron sâu (Neural/Deep Learning TTS)}: Tiêu biểu với các kiến trúc Tacotron 2, FastSpeech 2, hay VITS kết hợp neural vocoder (HiFi-GAN). Phương pháp này đem lại chất lượng ngữ điệu vượt bậc nhưng đòi hỏi kho ngữ liệu khổng lồ (hàng nghìn giờ âm thanh gán nhãn chi tiết) cùng năng lực tính toán phần cứng rất lớn để huấn luyện, không thích hợp cho các ứng dụng nhẹ triển khai tức thì trên thiết bị cục bộ.
\end{itemize}

Xuất phát từ yêu cầu học thuật của môn học Cơ sở dữ liệu đa phương tiện, nhóm nghiên cứu đã lựa chọn phương pháp \textbf{tổng hợp ghép nối mức từ (Word-level Concatenative TTS)} kết hợp với các kỹ thuật xử lý tín hiệu số (DSP) hiện đại. Đề tài tập trung vào việc tự xây dựng và quản trị kho ngữ liệu âm thanh thực nghiệm, thiết kế mô hình cơ sở dữ liệu quan hệ SQLite tối ưu hóa chỉ mục, đồng thời nghiên cứu giải thuật trích xuất các đặc trưng âm học số bao gồm: năng lượng RMS, tần suất cắt qua điểm không (ZCR), cao độ $F_0$ qua tự tương quan và 13 hệ số Cepstral trên thang tần số Mel (MFCCs).

Tính mới và giá trị ứng dụng của đề tài thể hiện ở việc:
\begin{enumerate}
    \item Xây dựng bộ quy ước mã hóa VNI hai chiều (Bidirectional VNI Encoding) độc lập với bảng mã Unicode, giải quyết triệt để sự phức tạp của ký tự tiếng Việt có dấu và cho phép tra cứu bản ghi đa phương tiện với tốc độ $O(1)$.
    \item Thiết kế và cài đặt thuật toán ghép nối mượt mà sử dụng kết hợp cắt tỉa khoảng lặng động (Silence Trimming), chuẩn hóa năng lượng RMS đồng nhất, kiểm soát biên độ đỉnh (Peak Limiter), mô hình hóa khoảng nghỉ ngắt câu (Pause Duration Modeling 120 ms) và kỹ thuật làm mờ chéo tuyến tính (Linear Cross-fade 10 ms).
    \item Tích hợp đầy đủ các API RESTful bất đồng bộ hiệu năng cao (FastAPI) cùng giao diện Web hiện đại (React 18 Single Page Application), hỗ trợ trực quan hóa dạng sóng, đường cong cao độ $F_0$ và phổ MFCC trực tiếp trên trình duyệt.
\end{enumerate}

Mục tiêu cốt lõi của đề tài là hoàn thiện hệ thống VietVoice TTS đáp ứng trọn vẹn 7 tiêu chí đánh giá khắt khe của học phần Cơ sở dữ liệu đa phương tiện: (1) Ngữ liệu song song văn bản -- âm thanh; (2) Kỹ thuật tổng hợp TTS mức từ; (3) Khả năng nhận văn bản và xuất âm thanh trực tiếp; (4) Sơ đồ khối và quy trình xử lý 5 bước rõ ràng; (5) Trích xuất và quản lý đặc trưng âm học số; (6) Mô hình CSDL đa phương tiện quan hệ chuẩn hóa; và (7) Thực nghiệm, đánh giá chất lượng toàn diện.

